% !TEX root = ../Thesis.tex

%---------------------------------------------------------------------------%
%->> Backmatter
%---------------------------------------------------------------------------%
\chapter[致谢]{致\quad 谢}\chaptermark{致\quad 谢}% syntax: \chapter[目录]{标题}\chaptermark{页眉}
%\thispagestyle{noheaderstyle}% 如果需要移除当前页的页眉
%\pagestyle{noheaderstyle}% 如果需要移除整章的页眉

此处填写致谢。


\rightline{2023年6月}
\chapter{作者简历及攻读学位期间发表的学术论文与其他相关学术成果}

\section*{作者简历：}
2013年09月——2017年06月，在武汉科技大学计算机学院获得学士学位。

2023年09月——2026年06月，在中国科学院大学杭州高等研究院物理与光电工程学院获得硕士学位。

\section*{已发表（或正式接受）的学术论文：}
{
\setlist[enumerate]{}% restore default behavior
\begin{enumerate}[nosep]
    \item \textbf{Xu J}, Ye Z, Wang S, et al. Dual holographic and polarization encoding for high fidelity image transmission through multimode fibers[J]. Optics \& Laser Technology, 2025, 191: 113301.
    \item Ye Z, \textbf{Xu J}, Wang S, et al. Computational immunity-belt reconstruction restores far-field images through bent step-index fibers[J]. Photonics Research, 2025, 13(11): 3199-3209.
    \item Huang S, Ye Z, \textbf{Xu J}, et al. Force Sensing in Multimode Fibers via Wavefront-Shaped Focal-Spot Analysis[J]. IEEE Sensors Journal, 2025.
\end{enumerate}
}

\section*{申请或已获得的专利：}
{
\setlist[enumerate]{}% restore default behavior
\begin{enumerate}[nosep]
    \item 基于多模光纤全息和偏振编码的图像重建方法及成像系统，申请号：202510705774.2
\end{enumerate}
}

\cleardoublepage[plain]% 让文档总是结束于偶数页，可根据需要设定页眉页脚样式，如 [noheaderstyle]
%---------------------------------------------------------------------------%
